\documentclass{article}
\usepackage[utf8]{inputenc}
\usepackage[T1]{fontenc}
\usepackage[a4paper,total={6in,8in}]{geometry}
\usepackage{amsmath,amssymb,amsfonts}
\usepackage{graphicx}
\usepackage{hyperref}
\usepackage{enumitem}
\usepackage{listings}
\usepackage{xcolor}
\DeclareMathOperator{\Var}{Var}
\DeclareMathOperator{\Cov}{Cov}
\providecommand{\tightlist}{\setlength{\itemsep}{0pt}\setlength{\parskip}{0pt}}
\lstset{basicstyle=\ttfamily\scriptsize,breaklines=true,columns=fullflexible,
  keepspaces=true,showstringspaces=false,frame=single}
\hypersetup{colorlinks=true,urlcolor=blue,linkcolor=blue}
\setlength{\emergencystretch}{2em}

\title{Math 126 Problem 16: Efficiently Computing the SVD}
\author{} % Put your name here.
\date{Spring 2027}
\begin{document}
\maketitle

\hypertarget{submission-instructions}{%
\subsection*{Submission Instructions}\label{submission-instructions}}

Please upload your work to Gradescope by \textbf{11:59pm on Monday, April 5, 2027 (Pacific time).}

\begin{itemize}
\item
  You are encouraged to discuss ideas and work with your classmates. However, you \emph{must acknowledge} your collaborators at the top of each solution on which you collaborated with others and you \emph{must write} your solutions independently.
\item
  Your solutions to theory questions must be written legibly, or typeset in LaTeX or markdown.
\item
  I recommend that you write your solutions in a Python notebook using Google Colab.
\item
  You should submit your solutions as a \emph{single PDF} via the assignment on Gradescope.
\end{itemize}

\textbf{Grading:} The point of the problem set is for \emph{you} to learn. To this end, I hope to disincentivize the use of LLMs by \emph{not} grading your work for correctness. Instead, you will grade your own work by comparing it to my solutions. This self-grade is due the Friday \emph{after} the problem set is due, also on Gradescope.

\noindent Original notebook: \href{https://colab.research.google.com/drive/10fug5ckPaH6QQvK16XiD74eFa-VBHqmW?usp=sharing}{Google Colab}.

% Source notebook cell 2
\hypertarget{efficiently-computing-the-svd}{%
\subsection*{Efficiently Computing the SVD}\label{efficiently-computing-the-svd}}

We will explore how to efficiently approximate the top singular vector of a data matrix \(\mathbf{X} \in \mathbb{R}^{n \times d}\). Instead of computing the full SVD (which takes \(O(nd^2)\) time), we will analyze and implement two iterative algorithms operating on the symmetric matrix \(\mathbf{A} = \mathbf{X}^\top \mathbf{X}\): the Power Method and the Lanczos Method.

\hypertarget{part-a-tridiagonal-recurrence}{%
\subsection*{Part A: Tridiagonal Recurrence}\label{part-a-tridiagonal-recurrence}}

The efficiency of the Lanczos method relies entirely on the fact that projecting a symmetric matrix \(\mathbf{A}\) onto its Krylov subspace results in a tridiagonal matrix. In this part, you will prove why this happens.

Let \(\mathbf{A} \in \mathbb{R}^{d \times d}\) be a symmetric matrix. Let \(\mathbf{Q} = [\mathbf{q}_1, \mathbf{q}_2, \dots, \mathbf{q}_q]\) be a matrix with orthonormal columns that span the degree \(q\) Krylov subspace generated by \(\mathbf{A}\) and an initial vector \(\mathbf{z}^{(0)}\). Let \(\mathbf{T} = \mathbf{Q}^\top \mathbf{A} \mathbf{Q}\).

\begin{enumerate}
\def\labelenumi{\arabic{enumi}.}
\item
  \textbf{The Hessenberg Property}: By the definition of the Krylov subspace, we know that \(\mathbf{A}\mathbf{q}_j\) is strictly contained within the span of the first \(j+1\) basis vectors. Use this fact to prove that \(T_{i,j} = 0\) for all \(i > j + 1\).
\item
  \textbf{Symmetry}: Prove that the matrix \(\mathbf{T}\) is symmetric (\(\mathbf{T} = \mathbf{T}^\top\)).
\end{enumerate}

Combine your results from steps 1 and 2 to rigorously conclude that \(T_{i,j} = 0\) for all \(|i - j| > 1\). What does this tell us about the structure of \(\mathbf{T}\)?

\hypertarget{part-b-power-method-algorithm}{%
\subsection*{Part B: Power Method Algorithm}\label{part-b-power-method-algorithm}}

Write a Python function to implement the Power Method. Your function should take a symmetric matrix \(\mathbf{A}\) and the number of iterations \(q\), and return the approximated top eigenvector.

\hypertarget{part-c-lanczos-algorithm}{%
\subsection*{Part C: Lanczos Algorithm}\label{part-c-lanczos-algorithm}}

Now, write a Python function to implement the Lanczos Method. Remember that in practice, we iteratively build \(\mathbf{Q}\) and \(\mathbf{T}\) using the three-term recurrence to ensure numerical stability and speed.

\textbf{Hint:} For finding the top eigenvector of your small \(q \times q\) matrix \(\mathbf{T}\), you may use \texttt{np.linalg.eigh}.

\hypertarget{part-d-comparison}{%
\subsection*{Part D: Comparison}\label{part-d-comparison}}

We will test these algorithms on the digits dataset from \texttt{scikit-learn}. Extract the data matrix \(\mathbf{X}\) and mean-center its columns to form the covariance matrix \(\mathbf{A} = \mathbf{X}^\top \mathbf{X}\).

Run both your power method and Lanczos method implementations on \(\mathbf{A}\) for \(q=20\) iterations. Plot the \(\ell_2\) distance between the estimated true top eigenvector by iteration. What do you notice?

\end{document}
